%% =================================================================
%% Chen et al. Measurement 264 (2026) 120300.
%% Reconstruction of page 9 (printed p. 9).
%% Contents: Fig. 12 (spherical-cap scanning-speed & step study,
%% 9 sub-panels), prose covering Sections 4.3 (accuracy) tail and
%% 4.4 (scanning speed & step).
%% =================================================================

\documentclass[11pt]{article}

\usepackage[margin=0.85in]{geometry}
\usepackage{amsmath, amssymb}
\usepackage{mathrsfs}
\usepackage{xcolor}
\usepackage{fancyhdr}

\pagestyle{fancy}
\fancyhf{}
\lhead{\small Z.~Chen et al.}
\rhead{\small Measurement 264 (2026) 120300}
\cfoot{\small 9 $\to$ \arabic{page}}
\renewcommand{\headrulewidth}{0.4pt}

\setcounter{page}{1}
\setcounter{equation}{1}
\setlength{\parskip}{6pt}
\setlength{\parindent}{0pt}

\begin{document}

% ---------- Fig. 12 placeholder ----------
\begin{center}
\fbox{\begin{minipage}[c][8cm][c]{0.95\textwidth}
\centering
\textbf{Fig. 12.}\ Scanning speed and scanning step of the tactile
tomography system.\\[4pt]
\textcolor{gray}{\footnotesize
\textbf{(a)--(c)} A spherical-cap model: (a) cross-section with
4 mm AB-glue layer atop ABS substrate; (b) 3D CAD view; (c)
photograph of the filled sample.\\[3pt]
\textbf{(d)--(i)} 3D imaging of compression-depth distribution
when the tactile tomography system scanned the spherical cap under
nine speed/step combinations:
\textbf{(d)} 0.6 mm/s \& 0.5 mm,
\textbf{(e)} 1.2 mm/s \& 0.5 mm,
\textbf{(f)} 1.8 mm/s \& 0.5 mm,
\textbf{(g)} 2.4 mm/s \& 0.5 mm,
\textbf{(h)} 2.4 mm/s \& 1.0 mm,
\textbf{(i)} 2.4 mm/s \& 2.0 mm.
Each panel renders a rainbow-colored 3D dome on a red substrate;
Z scale 0--3.6 mm. Panels (d)--(g) preserve the smooth cap; (h)
shows slight facet coarsening; (i) (2 mm step) visibly fails to
reproduce the cap shape.]}
\end{minipage}}
\end{center}

\vspace{0.4em}

\textbf{Fig. 12.}\ Scanning speed and scanning step of the tactile
tomography system. (a), (b) and (c) A spherical cap model filled
with a soft surface layer of AB glue. 3D imaging of compression
depth distribution when the tactile tomography system scanned the
spherical cap model with a soft surface layer under scanning speeds
and scanning steps of (d) 0.6 mm/s \& 0.5 mm, (e) 1.2 mm/s \& 0.5
mm, (f) 1.8 mm/s \& 0.5 mm, (g) 2.4 mm/s \& 0.5 mm, (h) 2.4 mm/s
\& 1.0 mm, and (i) 2.4 mm/s \& 2.0 mm, respectively.

\vspace{0.6em}

% ---------- Prose ----------
A gauge block with a length of 30 mm, a width of 9 mm, and a height
of 3 mm was filled with a soft surface layer of AB glue to
investigate the scanning accuracy of the tactile tomography system
(Fig.~11a). The scanning area is $20 \times 20$ mm$^2$ (Fig.~11b),
and the scanning step is 0.5 mm, resulting in a total number of
scanning points of 1600. As shown in Fig.~11c, the scanning points
corresponding to the surface of the gauge block have a relatively
uniform compression depth. The compression depth of the scanning
points in other regions presents a slight difference due to the
incomplete smoothing of the bottom, which was prepared by 3D
printing with ABS. As a result, a scanning accuracy of over 98.9\%
is distributed throughout the scanning area (Fig.~11d), and the
scanning area corresponding to the surface of the gauge block shows
a high scanning accuracy of over 99.3\%, demonstrating that the
tactile tomography system has the ability to scan every point
accurately and stably.

\subsection*{4.4.\ Scanning speed and scanning step of the tactile tomography system}

There is a trade-off between the scanning efficiency and imaging
quality for scanning imaging, in which the scanning efficiency is
determined by the scanning speed and scanning step. As shown in
Fig.~12a, 12b, and 12c, a spherical cap model was filled with a
soft surface layer of AB glue to investigate the scanning
efficiency and imaging quality of the tactile tomography system.
The spherical cap model was scanned by the tactile tomography
system with scanning speeds and scanning steps of 0.6 mm/s \& 0.5
mm, 1.2 mm/s \& 0.5 mm, 1.8 mm/s \& 0.5 mm, 2.4 mm/s \& 0.5 mm,
2.4 mm/s \& 1.0 mm, and 2.4 mm/s \& 2.0 mm, respectively. The
reconstructed 3D images corresponding to the scanning speeds of
0.6, 1.2, 1.8, and 2.4 mm/s are shown in Fig.~12d, e, f, and g,
respectively, and all these 3D images can reproduce the spherical
cap model. The curved shape of the 3D images remained highly smooth
as the scanning speed increased to 2.4 mm/s, demonstrating that
the increase of scanning speed will not greatly degrade the imaging
quality. Fig.~12g, h, and i show the reconstructed 3D images
corresponding to the scanning speed and scanning step of 2.4 mm/s
\& 0.5 mm, 2.4 mm/s \& 1.0 mm, and 2.4 mm/s \& 2.0 mm,
respectively. There is a relative decrease in the smoothness of
the curved shape of the 3D images as the scanning step increased,
and the 3D image could even not reproduce the spherical cap model
well when the scanning step increased to 2 mm. The

\end{document}
